% ---------------------------------------------------------------------------
% Section: Mathematical formulation
% ---------------------------------------------------------------------------
\section{Mathematical Formulation}
\label{sec:math}

This section turns the inspiration of Section~\ref{sec:background} into a formal
model. Nothing here depends on the metaphysical reading: the objects are
policies, verifiers, costs and losses. Each result ends with its
\emph{engineering consequence}, which the specification (Section~\ref{sec:spec})
turns into a parameter or an algorithm. Results are labelled by the strength of
their assumptions: \textbf{[exact]} holds under the stated model without further
hypotheses; \textbf{[strong]} needs an assumption that must be checked
empirically; \textbf{[conjecture]} is a design hypothesis scheduled for
validation in the roadmap.

\subsection{Objects}
\label{sec:objects}

\begin{definition}[Task, answer, correctness]
A \emph{task} is $x\in\mathcal X$ (a query together with its conversational
context). An \emph{answer} is $y\in\mathcal Y$. Correctness is an unobserved
indicator $\kappa(x,y)\in\{0,1\}$. A task is \emph{verifiable} if a deterministic
procedure can evaluate $\kappa$ (unit tests, numeric or symbolic equality,
schema validity); otherwise it is \emph{non-verifiable}.
\end{definition}

\begin{definition}[Fast and slow paths]
The \emph{fast path} (System~1, the analogue of the impulsive
\emph{nafs}) is a stochastic policy $\pi_1(y\mid x)$ with expected cost
$c_1(x)$ and correctness probability $p_1(x)=\Pr_{y\sim\pi_1}[\kappa(x,y)=1]$.
The \emph{slow path} (System~2, deliberation) is a policy $\pi_2(y\mid x,b)$
parametrised by a \emph{thinking budget} $b\in\mathbb N$ (maximum reasoning
tokens), with cost $c_2(x,b)$ and correctness $p_2(x,b)$. Costs are measured in
a common unit (seconds of CPU time or tokens weighted by a tariff).
\end{definition}

\begin{definition}[Verifier]
A \emph{verifier} is a map $V:\mathcal X\times\mathcal Y\to\{0,1\}$ with
true-accept rate $r=\Pr[V=1\mid\kappa=1]$ and false-accept rate
$\varphi=\Pr[V=1\mid\kappa=0]$. A verifier is \emph{exact} if $\varphi=0$ and
$r=1$ on its domain. Soft verifiers (rubric judges, consensus) have
$\varphi>0$ and must publish measured $(r,\varphi)$.
\end{definition}

\begin{definition}[Gate and trace state]
After the fast path has run, the \emph{controller} (the \emph{\textquoteleft aql}
analogue) observes a feature vector $s=\sigma(x,\text{trace}_1)\in\mathcal S$
(token log-probabilities, agreement among cheap samples, verifier outcome,
domain class) and chooses, with probability $\gamma(s)\in[0,1]$, to
\emph{escalate} to the slow path. The lifecycle of one request is a finite
automaton over the \emph{trace states}
$\mathcal Z=\{\mathsf{A},\mathsf{L},\mathsf{M},\mathsf{T},\mathsf{X}\}$ of
Fig.~\ref{fig:nafs-automaton}: $\mathsf A$ (an unverified draft),
$\mathsf L$ (the draft has been flagged by the gate or a verifier),
$\mathsf M$ (a deliberate retry conditioned on evidence), $\mathsf T$ (an
accepted answer) and $\mathsf X$ (abstention with an explicit uncertainty
notice). The labels are mnemonic for the four classical stages of the
\emph{nafs} (Section~\ref{sec:islamic}); the automaton itself is an ordinary
bounded retry loop.
\end{definition}

\begin{definition}[Losses and objective]
A wrong answer costs $L>0$, an abstention costs $L_a\in(0,L)$, and compute is
priced by $\lambda>0$ per cost unit. For a controller
$\theta=(\gamma,b,N,R)$ (gate, thinking budget, samples per round, retry
rounds) the design objective is
\begin{equation}
\label{eq:objective}
\begin{split}
\mathcal J(\theta)=\mathbb E_x\big[&L\,\Pr[\text{wrong}\mid x,\theta]
+L_a\,\Pr[\text{abstain}\mid x,\theta]\\
&+\lambda\,\mathbb E[\text{cost}\mid x,\theta]\big].
\end{split}
\end{equation}
\end{definition}

\begin{lemma}[Termination] \textnormal{[exact]}
\label{lem:termination}
With at most $N$ slow-path samples per round and at most $R$ retry rounds, every
request makes at most $1+N(R{+}1)$ model calls and at most $1+N(R{+}1)$
verifier calls before reaching $\mathsf T$ or $\mathsf X$.
\end{lemma}
\begin{proof}
Each transition of the automaton either consumes one of the $N(R{+}1)$ slow
samples or terminates; the fast path is called once. \qed
\end{proof}
\emph{Consequence.} Every request has a hard upper bound on latency
(NFR-001 in the specification); the gateway MUST enforce $N$ and $R$ as
configuration limits rather than letting the model decide when to stop.

\subsection{Precision of a verified answer}

\begin{theorem}[Verified-return precision] \textnormal{[exact]}
\label{thm:precision}
Let a candidate be correct with prior probability $p$ and let it be accepted
by a verifier with rates $(r,\varphi)$. Then
\begin{equation}
\rho(p)\;=\;\Pr[\text{wrong}\mid\text{accepted}]
=\frac{\varphi(1-p)}{\varphi(1-p)+r\,p},
\end{equation}
and $\rho(p)\le\varepsilon$ if and only if
$\varphi\le\dfrac{\varepsilon\,r\,p}{(1-\varepsilon)(1-p)}$. If $k$ verifiers
are applied and their acceptances are conditionally independent given
$\kappa$, the same holds with $\varphi\leftarrow\prod_i\varphi_i$ and
$r\leftarrow\prod_i r_i$.
\end{theorem}
\begin{proof}
Bayes' rule with $\Pr[\text{accepted}]=\varphi(1-p)+rp$. Solving
$\rho\le\varepsilon$ for $\varphi$ gives the bound. Conditional independence
factorises the joint acceptance probabilities. \qed
\end{proof}
\emph{Consequence.} The precision of the system on verifiable tasks is
bounded by the verifier's false-accept rate, not by the strength of either
model. Every verifier in the hub MUST publish $(r,\varphi)$ measured by
mutation testing on a held-out set (FR-030), and the per-verifier acceptance
threshold is $1.0$ for exact verifiers. For hard tasks ($p=0.1$) a $1\%$
precision target requires $\varphi\le1.1\times10^{-3}$: only exact verifiers
qualify, which is why consensus (Theorem~\ref{thm:plurality}) is a flag and not
a verifier.

\subsection{The gate is a threshold rule}

Let $e_1(s)=\Pr[\text{fast answer wrong}\mid s]$,
$e_2(s)=\Pr[\text{final answer wrong after escalation}\mid s]$ and
$\Delta c(s)=\mathbb E[\text{extra cost of escalation}\mid s]$.

\begin{theorem}[Optimal gate] \textnormal{[exact]}
\label{thm:gate}
\begin{enumerate}[label=(\roman*)]
\item The gate minimising \eqref{eq:objective} escalates if and only if
\begin{equation}
\label{eq:gate}
L\,\big(e_1(s)-e_2(s)\big)\;>\;\lambda\,\Delta c(s),
\end{equation}
with ties broken arbitrarily.
\item Under a budget constraint $\mathbb E[\gamma(s)\Delta c(s)]\le B$
instead of the price $\lambda$, the optimal gate escalates if and only if
$\big(e_1(s)-e_2(s)\big)/\Delta c(s)>\lambda^\star$, where $\lambda^\star$ is
the smallest multiplier that meets the budget; randomisation is needed only on
the level set $\{s:(e_1-e_2)/\Delta c=\lambda^\star\}$.
\item If $e_1-e_2$ is non-decreasing in a scalar score $s$ and $\Delta c$ is
constant, the rule is a single threshold $s>\tau$.
\end{enumerate}
\end{theorem}
\begin{proof}
(i) Conditional on $s$, keeping costs $Le_1(s)$ and escalating costs
$\lambda\Delta c(s)+Le_2(s)$; the objective is an expectation over $s$ of a
pointwise choice, so pointwise minimisation is optimal. (ii) The Lagrangian
$\mathbb E[\gamma(s)(L e_2+\lambda\Delta c-Le_1)]$ is linear in $\gamma$; by the
Neyman--Pearson lemma the optimal $\gamma$ is an indicator of a likelihood-ratio
set with randomisation on the boundary. (iii) Immediate from (i). \qed
\end{proof}
\emph{Consequence.} The controller is implemented as a calibrated estimator
$\hat e_1(s)$ (logistic or isotonic regression on the fast-path features),
a per-domain estimate $\hat e_2$, and the comparison \eqref{eq:gate}. The only
tunable is the cost ratio $\lambda\Delta c/L$ (or equivalently the budget $B$),
which the operator sets; the threshold is \emph{derived}, not hand-tuned
(FR-010--FR-013). The rule coincides with Chow's optimal reject rule
\cite{chow1970reject} and with the expected-value-of-control account of
cognitive control \cite{shenhav2013evc}, with ``reject'' read as
``deliberate''.

\begin{proposition}[Price of miscalibration] \textnormal{[exact]}
\label{prop:miscal}
If the plug-in gate uses estimates with
$\sup_s|\hat e_1(s)-e_1(s)|\le\epsilon_1$ and
$\sup_s|\hat e_2(s)-e_2(s)|\le\epsilon_2$, its objective exceeds that of the
optimal gate by at most $L(\epsilon_1+\epsilon_2)$.
\end{proposition}
\begin{proof}
At any $s$ the two gates disagree only if
$|L(e_1-e_2)-\lambda\Delta c|\le L(\epsilon_1+\epsilon_2)$, and the loss from a
wrong decision equals that absolute difference. Taking expectations over $s$
gives the bound. \qed
\end{proof}
\emph{Consequence.} The calibration target of the gate (NFR-020) is set from
the regret the operator tolerates: a maximum calibration error of $0.05$ on
$\hat e_1$ bounds the excess loss by $0.05L$ per request.

\begin{proposition}[Cost--quality frontier] \textnormal{[exact]}
\label{prop:frontier}
The set of achievable pairs (expected cost, expected error) over all gates
$\gamma$ is convex, and its lower-left boundary is traced by the family of
threshold gates of Theorem~\ref{thm:gate} as $\lambda$ varies.
\end{proposition}
\begin{proof}
Both coordinates are linear in $\gamma$, so mixtures of gates achieve
mixtures of pairs; the boundary of a convex set is supported by linear
functionals, each of which is minimised by a threshold gate by
Theorem~\ref{thm:gate}. \qed
\end{proof}
\emph{Consequence.} The user-facing ``speed--care'' control in the chat UI is
a single scalar $\lambda$ (FR-014); no other knob can reach a better point.

\subsection{Deliberation: sampling, stopping and retries}

Let a slow-path round draw one candidate and verify it; let
$a=q\,r+(1-q)\varphi$ be the probability that a draw is accepted, where $q$ is
the probability that a draw is correct.

\begin{theorem}[Coverage and cost of sequential verified sampling]
\textnormal{[exact]}
\label{thm:coverage}
With at most $M$ draws and stopping at the first acceptance,
(i) the probability that some candidate is accepted is $1-(1-a)^M$;
(ii) the expected number of draws is $E(M)=\big(1-(1-a)^M\big)/a$;
(iii) the conditional error of the accepted candidate is $\rho(q)$ of
Theorem~\ref{thm:precision}, independent of $M$;
(iv) if $a$ varies across tasks with mean $\bar a$, the coverage is at most
$1-(1-\bar a)^M$.
\end{theorem}
\begin{proof}
(i)--(ii) are the truncated geometric distribution. (iii) Each accepted
candidate, whichever draw produced it, is wrong with probability $\rho(q)$ by
Bayes' rule. (iv) $u\mapsto(1-u)^M$ is convex, so Jensen's inequality gives
$\mathbb E[(1-a)^M]\ge(1-\bar a)^M$. \qed
\end{proof}
\emph{Consequence.} More samples buy \emph{coverage}, never \emph{precision}:
the number of draws is chosen from a coverage target and the per-task
heterogeneity of $a$ must be measured, not averaged (Phase~0 of the roadmap).

\begin{proposition}[Optimal number of draws under uncertainty about $a$]
\textnormal{[strong: Beta prior, monotone stopping]}
\label{prop:stopping}
Suppose the acceptance probability of a task is unknown with prior
$a\sim\mathrm{Beta}(\alpha,\beta)$, each draw costs $c_p$, and failing to
produce an accepted candidate costs $L_a$. After $m$ consecutive rejections the
posterior mean is $a_m=\alpha/(\alpha+\beta+m)$. The policy that draws again if
and only if $a_m L_a>c_p$ is optimal, so the optimal cap is
\begin{equation}
M^\star=\Big\lceil\frac{\alpha L_a}{c_p}-\alpha-\beta\Big\rceil
\end{equation}
(the number of $m\ge0$ with $\alpha/(\alpha+\beta+m)>c_p/L_a$),
when the right-hand side is positive, and $M^\star=0$ otherwise.
\end{proposition}
\begin{proof}
A further draw costs $c_p$ and, with posterior-predictive probability $a_m$,
saves $L_a$; so the one-step look-ahead rule continues iff $a_mL_a>c_p$.
Because every continuation happens only after a rejection and $a_m$ is
decreasing in $m$, once the rule says ``stop'' it says ``stop'' in all
subsequent states: the problem is monotone, and in the monotone case the
one-step look-ahead rule is optimal \cite{chow1971great}. The draws made are exactly those indexed by the $m\ge0$ with
$m<\alpha L_a/c_p-\alpha-\beta$, whose count is the ceiling. \qed
\end{proof}
\emph{Consequence.} The retry limit $R$ and the samples per round $N$
(with $M=N(R{+}1)$) are derived per domain from the Phase-0 fit of
$(\alpha,\beta)$ and the operator's cost ratio $c_p/L_a$ (FR-021). The Beta
model is a modelling choice; the roadmap validates it against the empirical
distribution of acceptance rates (hypothesis H3).

\subsection{Self-monitoring: why the loop must be verifier-gated}

The classical \emph{muhasaba} (self-accounting) is a loop in which a draft is
re-examined and amended. Let $u$ be the probability that one unverified
refinement turns a wrong answer into a right one and $d$ the probability that
it turns a right answer into a wrong one.

\begin{theorem}[Unverified refinement converges to a fixed point]
\textnormal{[strong: homogeneous Markov refinement]}
\label{thm:selfrefine}
If $q_t$ is the probability that the $t$-th draft is correct and
$q_{t+1}=q_t(1-d)+(1-q_t)u$, then
$q_t=q^\star+(q_0-q^\star)(1-u-d)^t$ with $q^\star=u/(u+d)$. In particular
$q_t$ decreases monotonically in $t$ whenever $q_0>q^\star$.
\end{theorem}
\begin{proof}
The recursion is affine with fixed point $q^\star$ and contraction factor
$1-u-d\in(-1,1)$ for $u+d\in(0,2)$. \qed
\end{proof}
\emph{Consequence.} Iterating self-critique without external evidence can only
help tasks whose first draft is worse than the fixed point; for a strong first
draft it \emph{harms}. This is consistent with the empirical finding that
language models do not reliably self-correct without external feedback
\cite{huang2024selfcorrect,kamoi2024survey}. Nafs-Aql therefore never accepts an
unverified refinement: on non-verifiable tasks the loop is capped at $R=1$
(FR-022).

\begin{theorem}[Verifier-gated loop is stable] \textnormal{[exact]}
\label{thm:gatedloop}
In a loop that accepts a candidate only when a verifier with rates
$(r,\varphi)$ passes it, and whose $t$-th candidate is correct with
probability $q_t$, the probability of terminating with an accepted wrong answer
is at most $\varphi\,\mathbb E[T]$, where $T$ is the number of candidates
examined, and the error conditional on acceptance lies in
$[\min_t\rho(q_t),\max_t\rho(q_t)]$. Neither bound depends on the refinement
dynamics $(u,d)$.
\end{theorem}
\begin{proof}
A wrong candidate at round $t$ is accepted with probability $\varphi$;
summing $\Pr[\text{reach }t]\,(1-q_t)\varphi\le\varphi\Pr[\text{reach }t]$
over $t$ gives $\varphi\,\mathbb E[T]$. Conditional on acceptance at round $t$
the error is $\rho(q_t)$ by Theorem~\ref{thm:precision}; the overall
conditional error is a mixture of these values. \qed
\end{proof}
\emph{Consequence.} With a verifier in the loop the risk is linear in
$\varphi$ and in the expected number of rounds, both of which the controller
bounds (Lemma~\ref{lem:termination}); the refinement prompt can be changed
freely without re-deriving the safety case.

\subsection{Consensus has a floor}

\begin{theorem}[Plurality voting converges to the modal answer]
\textnormal{[exact]}
\label{thm:plurality}
Let $N$ answers be drawn i.i.d.\ from $\pi(\cdot\mid x)$ over a finite answer
set with a unique mode $y^\circ(x)$. Then the plurality vote converges almost
surely to $y^\circ(x)$ as $N\to\infty$, so the limiting accuracy of
self-consistency is $\Pr_x[\,y^\circ(x)\text{ is correct}\,]$. For finite $N$,
if the correct answer has mass $q$ and the strongest wrong answer has mass
$\bar q<q$, the probability that the wrong answer ties or beats the correct one
is at most $\exp\!\big(-N(q-\bar q)^2/2\big)$.
\end{theorem}
\begin{proof}
The empirical frequencies converge almost surely to $\pi(\cdot\mid x)$
(strong law), so the arg-max converges to the unique mode. For the finite
bound let $D_i=\mathbb 1[y_i=y^\star]-\mathbb 1[y_i=\bar y]\in[-1,1]$ with mean
$q-\bar q$; Hoeffding's inequality \cite{hoeffding1963} gives
$\Pr[\sum_iD_i\le0]\le\exp(-2N(q-\bar q)^2/4)$. \qed
\end{proof}
\emph{Consequence.} Agreement among samples cannot repair tasks on which the
model is confidently wrong; the mass of such tasks is a floor that no $N$
removes. Consensus is therefore recorded as an \emph{assurance level}
(\texttt{consensus}) distinct from \texttt{verified} (FR-031), is never used
to write training data without a human check, and the number of samples for a
consensus flag is set from the margin $q-\bar q$ measured in Phase~0.

\subsection{Independence of the two paths}

\begin{proposition}[Anchoring and collision] \textnormal{[exact]}
\label{prop:collision}
If the slow path does not see the fast draft, the probability that both
produce the same wrong answer is the collision probability
$\sum_{y:\kappa(x,y)=0}\pi_1(y\mid x)\pi_2(y\mid x)$. Passing the draft to the
slow path can only add anchoring on top of this value. Certifying an adoption
rate below $5\%$ at confidence $95\%$ with zero observed adoptions needs at
least $\lceil\ln0.05/\ln0.95\rceil=59$ trials; with observed adoptions the
Clopper--Pearson interval \cite{clopper1934} applies.
\end{proposition}
\emph{Consequence.} The slow path is \emph{draft-blind} by default
(FR-023); draft-conditioned deliberation is an ablation, and the coupling test
(hypothesis H5) is a release gate. Because both paths may share weights (the
single-model profile of Section~\ref{sec:models}), the collision term is not
small a priori and is measured.

\subsection{Honest self-knowledge under routing}

\begin{proposition}[Unbiased capability estimates need audits]
\textnormal{[exact]}
\label{prop:audit}
Slow-path correctness is observed only on escalated tasks, so the naive
estimate is $\Pr[\text{slow correct}\mid\text{escalated}]$, a biased estimate of
$p_2$. If a random fraction $\epsilon>0$ of the kept tasks is also run through
the slow path, then with escalation rate $g$,
$\hat p_2=g\,\hat p_{2\mid\mathrm{esc}}+(1-g)\,\hat p_{2\mid\mathrm{audit}}$ is
unbiased, with standard error at most
$\tfrac{1}{2}\sqrt{\big(g+(1-g)/\epsilon\big)/n}$ over $n$ requests.
\end{proposition}
\begin{proof}
Inverse-probability weighting with known strata weights $g$ and $1-g$ and
sampling probabilities $1$ and $\epsilon$; the variance of each stratum mean is
at most $1/(4\,n_{\text{stratum}})$ with $n_{\text{esc}}=ng$ and
$n_{\text{audit}}=n(1-g)\epsilon$. \qed
\end{proof}
\emph{Consequence.} The self-model (the \emph{muhasaba} ledger) is built only
from traces plus a $5\%$ audit stream (FR-040); capabilities with no audited
data are reported as \emph{unmeasured}, never as zero. With $n=10^4$, $g=0.3$
and $\epsilon=0.05$ the standard error is at most $0.019$.

\subsection{Estimation and gate sample sizes}

\begin{proposition}[Sample sizes] \textnormal{[exact]}
\label{prop:samples}
\begin{enumerate}[label=(\roman*)]
\item (Threshold from a quantile.) If $\tau$ is set as the empirical
$(1-g^\star)$-quantile of $n$ i.i.d.\ scores, the realised escalation rate is
within $\varepsilon$ of $g^\star$ with probability at least $1-\delta$ whenever
$n\ge\ln(2/\delta)/(2\varepsilon^2)$ \cite{massart1990dkw}.
\item (Non-inferiority.) To show that the cascade loses at most $m$ points of
accuracy against the slow-only baseline with power $1-\beta$ at level $\alpha$,
a paired design with discordance rate $\delta$ needs
$n\ge(z_\alpha+z_\beta)^2\,\delta/m^2$ pairs.
\item (Binomial gates.) Certifying a rate below $\varepsilon$ with zero failures
at confidence $1-\delta$ needs $n\ge\ln\delta/\ln(1-\varepsilon)$ trials.
\end{enumerate}
\end{proposition}
\emph{Consequence.} The Phase-0 baseline is at least $4\,600$ tasks
($\varepsilon=0.02$, $\delta=0.05$), the one-point non-inferiority gate needs
about $6\,200$ pairs at $\delta=0.1$ (or $700$ for a three-point margin), and
every release gate in the roadmap states its sample size. With $n=500$ the
escalation rate is known only to $\pm6$ points.

\subsection{Distillation (\emph{tazkiya}) without self-deception}

\begin{proposition}[Label noise and verifier-inflated gains] \textnormal{[exact]}
\label{prop:distill}
\begin{enumerate}[label=(\roman*)]
\item If a fraction $\eta$ of training labels is wrong, then for every
hypothesis $h$, $|\mathrm{err}_{\text{noisy}}(h)-\mathrm{err}_{\text{clean}}(h)|\le\eta$;
with verifier-filtered traces $\eta=\rho$ of Theorem~\ref{thm:precision}.
\item If accuracy is measured with a verifier of rates $(r,\varphi)$, the
measured rate is $\varphi+(r-\varphi)\,t$ where $t$ is the true rate; a measured
gain $\Delta$ corresponds to a true gain $\Delta/(r-\varphi)$, and training on
verifier-accepted outputs can raise $\varphi$ and inflate the measurement.
\item A paired promotion gate with discordance $\delta$ resolving a gain of
$\varepsilon$ needs $n\ge z^2\delta/\varepsilon^2$ pairs per domain.
\end{enumerate}
\end{proposition}
\begin{proof}
(i) The two error events differ only on examples whose label is wrong.
(ii) Expand $\Pr[V=1]=r\,t+\varphi(1-t)$. (iii) Normal approximation to the
paired binomial. \qed
\end{proof}
\emph{Consequence.} Promotion of a distilled fast-path adapter requires a
confidence-bound gain on a verifier that was never used for data collection,
plus a per-domain non-regression bound and human sign-off (FR-050--FR-053).

\subsection{What a CPU can serve}

\begin{proposition}[Decode and prefill ceilings] \textnormal{[exact given the
roofline model]}
\label{prop:roofline}
For a dense decoder whose active weights occupy $W$ bytes and whose key--value
cache at context length $\ell$ occupies $K(\ell)$ bytes, the steady-state
decode rate on a machine with memory bandwidth $\beta$ is at most
$\beta/(W+K(\ell))$ tokens per second, and the prefill rate on a machine with
peak throughput $F$ FLOP/s and $P$ active parameters is at most $F/(2P)$
tokens per second.
\end{proposition}
\begin{proof}
Each decoded token must stream every active weight and the cache from memory
at least once; each prefill token requires at least $2P$ floating-point
operations. \qed
\end{proof}
\emph{Consequence.} On the reference machine (dual-channel DDR4-2400,
$\beta\approx38$\,GB/s; AVX2, four cores) a 4-billion-parameter model in 4-bit
GGUF ($W\approx2.5$\,GB) cannot exceed about 15 tokens/s, a 1.7B model about
35 tokens/s, and a 1000-token prefill of the 4B model takes seconds to tens of
seconds. These ceilings fix the model sizes, the thinking-budget cap, the
context cap and the prompt-cache requirement (NFR-002--NFR-006). Measured
values on the reference machine are reported in Section~\ref{sec:feasibility}
and are below the ceilings, as expected.

\begin{proposition}[Thinking budget] \textnormal{[strong: concave returns]}
\label{prop:budget}
If $p_2(x,b)$ is concave and non-decreasing in $b$ and the cost is linear,
$c_2(x,b)=c_2^0+\kappa_{\text{tok}}\,b$, the budget minimising
$L(1-p_2(x,b))+\lambda c_2(x,b)$ is the unique $b^\star$ with
$L\,\partial_bp_2(x,b^\star)=\lambda\kappa_{\text{tok}}$ (or the boundary if
no interior solution exists).
\end{proposition}
\emph{Consequence.} The thinking budget is a per-domain configuration value
estimated from the Phase-0 accuracy-versus-budget curve (hypothesis H4);
when the budget is exhausted the controller forces the answer by closing the
reasoning span \cite{muennighoff2025s1} rather than returning nothing.

\subsection{End-to-end identity}

\begin{proposition}[Cascade difference] \textnormal{[exact]}
\label{prop:cascade}
Let $\rho_1$ be the error on tasks the gate keeps and $\bar e_2^{\,\text{kept}}$
the error the slow path alone would make on those tasks. If the slow path
behaves identically on escalated tasks in both systems,
$\Pr[\text{cascade wrong}]-\Pr[\text{slow-only wrong}]
=(1-g)\big(\rho_1-\bar e_2^{\,\text{kept}}\big)$.
\end{proposition}
\emph{Consequence.} Non-inferiority within $m$ requires
$\rho_1\le\bar e_2^{\,\text{kept}}+m/(1-g)$; $\bar e_2^{\,\text{kept}}$ is
observable only through the audit stream of Proposition~\ref{prop:audit}, which
is therefore part of the safety case.

\subsection{Summary of engineering consequences}
Table~\ref{tab:consequences} lists each result with the parameter or
algorithm it determines; the specification's traceability matrix repeats the
mapping by requirement ID.

\begin{table}[t]
\caption{Formal results and what they determine}
\label{tab:consequences}
\centering\small
\begin{tabular}{@{}p{2.2cm}p{5.8cm}@{}}
\toprule
Result & Engineering consequence \\
\midrule
Lemma~\ref{lem:termination} & Hard caps $N$, $R$; latency bound \\
Thm.~\ref{thm:precision} & Verifier $(r,\varphi)$ publication; exact-verifier threshold $1.0$ \\
Thm.~\ref{thm:gate} & Gate $=$ calibrated $\hat e_1$, $\hat e_2$, cost ratio; no hand-tuned $\tau$ \\
Prop.~\ref{prop:miscal} & Calibration-error target of the gate \\
Prop.~\ref{prop:frontier} & Single ``speed--care'' scalar in the UI \\
Thm.~\ref{thm:coverage} & $N$ from coverage target; measure heterogeneity of $a$ \\
Prop.~\ref{prop:stopping} & $R$, $N$ from Beta fit and $c_p/L_a$ \\
Thm.~\ref{thm:selfrefine} & No unverified refinement; $R=1$ on non-verifiable tasks \\
Thm.~\ref{thm:gatedloop} & Loop risk $\le\varphi\,\mathbb E[T]$; prompt changes need no new safety case \\
Thm.~\ref{thm:plurality} & Consensus is an assurance flag; $N$ from margin \\
Prop.~\ref{prop:collision} & Draft-blind slow path; coupling test gate \\
Prop.~\ref{prop:audit} & $5\%$ audit stream; ``unmeasured'' label \\
Prop.~\ref{prop:samples} & Phase-0 size $\ge4\,600$; gate sample sizes \\
Prop.~\ref{prop:distill} & Promotion gate: held-out verifier, confidence bounds \\
Prop.~\ref{prop:roofline} & Model sizes, context cap, prompt cache, budget cap \\
Prop.~\ref{prop:budget} & Per-domain thinking budget; budget forcing \\
Prop.~\ref{prop:cascade} & Audit stream is part of the safety case \\
\bottomrule
\end{tabular}
\end{table}
